\RequirePackage{plautopatch}
\documentclass[dvipdfmx]{jarticle}
\DeclareUnicodeCharacter{200B}{}
\usepackage{mabst}

\title{
情報システム工学専攻 中間発表会予稿 予稿
}

\newcommand{\ID}{25MM356}
\newcommand{\Name}{LIU RUNZHI}
\newcommand{\Teacher}{小林 貴訓}

%\setlength{\oddsidemargin}{-2mm}  % プリンターに合わせて調整して下さい

\begin{document}
\maketitle

\section{はじめに}
 近年の少子高齢化に伴い，外出が難しい高齢者が増え，買い物に付随する会話や気分転換といった社会的機会が失われつつある．
 買い物は単なる購買にとどまらず健康維持に資する活動として日常的に取り入れられている一方，人手不足により十分な同伴支援が困難であることが課題である．
 この課題に対して，遠隔地からの買い物参加を可能にする通信基盤や，ユーザの対話指示に基づき売場内を自律走行する買い物カートが開発されてきた．
 
 しかし、現在の買い物カートは，決まったコマンドや定型のやり取りに頼ることが多く，実際の売り場では何度も確認する会話や，
 目標の細かい指定，利用者への提案がうまくできず，現地ユーザーとの双方向のやり取りが弱いという課題がある．
 そこで本研究では大規模言語モデル（LLM）を用い，遠隔利用者の自然な発話を実行可能な「行き先（エリア）」とタスクの並びに変換する．
 ロボットがエリアに到着した後LLMは，人物と目標物の位置を合わせて観測し，
 F-formationなどの社会的モデルに基づいて，人に向き，かつ目標物と整合する立ち位置と停止条件（方位・距離の閾値と保持時間）を生成する．
 これにより，カメラの遮蔽を減らし，遠隔と現地で同じ対象を見る「共同注意」を高め，自然なやり取りを実現する．

\section{ユーザに追従する自律移動手法}
  テレプレゼンスロボット(telepresence robot)においては，遠隔操作者が移動制御に注力すると認知負荷が高まり，会話の質・安全性・自然さが低下しやすい．
したがって，ロボット側が自律的に人物を追従し，障害物回避と快適な相対距離・方位を維持する機能が求められる．
Cosgun らは，この要請に応えるべく,社会的に望ましい自主人物追従を実現した [3]．
提案手法では，追従対象となる人の将来軌道を見越し，時間とともに変化する目標関数とコスト関数の重み付き和に対する割引付き効用を最大化する経路を探索する．
コスト関数には障害物との距離コスト，加速度コスト，角速度コストを導入し，障害物回避・急加減速の抑制・蛇行の抑制を同時に満たす軌道を選好させる．
一方，ゴール（目標）関数はロボットと人の相対位置関係に基づく評価を与え，対話しやすい人の右斜め後方付近に位置し，社会的に妥当な距離帯を保つ配置を促す．
．ユーザスタディの結果，遠隔操作者による手動操作に比べ，社会的空間を考慮した自律追従は安全性・使いやすさ・対話の自然さを有意に向上させ，
遠隔ユーザがロボット制御ではなく会話に集中できることが示された．
  つまり、社会ルールに適合した移動ロボットを作成するには、人を追跡する機能が不可欠と考えられる．
  
\section{グループ認識に基づく自律移動手法}
  Malladiらは，混雑環境でのロボット移動が個人の回避でグループの間を割る等を招きやすいという背景から，
  周囲の人をDBSCAN(Den-sity Based Spatial Clustering of Applications with Noise)で動的にクラスタ化し「グループ」を一級の回避対象として扱う社会的ナビ手法 SANGO を提案した[4]．
  ロボットは「自分の位置、壁などの静的障害物、人やグループとの相対的な距離と向き」を入力として受け取り、
  グループに近づきすぎるほど減点になる点数ルールのもとで、PPOを使って最適な動き方を学習する．
  実験は，自作の2Dシミュレータ（MOSANG/COG）で，人数密度や配置を変えた複数シナリオを用い，成功率・衝突率・最小衝突時間・グループ侵入度（不快度）・停滞時間などを評価指標として比較した．
  結果として，従来の“個人ベース回避”に比べて不快度と衝突率を大きく低減し（複雑環境で大幅減），停滞時間を短縮しながら到達成功率を有意に向上させ，
  グループ間を社会的に妥当な距離で通過できることを示した．社会的に望ましい移動ロボットを実現するためには，個人単位の回避にとどまらず，群集内のグループ関係を考慮することが不可欠である
\section{コミュニケーションを考慮した自律移動手法}
 ロボットが人と自然に協力するためには，相手に分かりやすい振る舞いで意図を可視化し，予測しやすく安全・安心な行動をとることが重要である．［5］
 この目標を達成するには，ロボットが人とより自然にコミュニケーションできる仕組みを備えることが鍵となる．


\subsection{人間とロボットの共同注意}
 共同注意とは，人とロボットが同じ対象に注意を向け，その共有が互いに確認されている状態を指す．
 Huangらは机上タスクでロボットが人の視線・指さしに応答する条件としない条件を比較し，
 動画評価でロボットが注意を監視し・確保（ずれたら見返しや再指示で戻す）する振る舞いの有無を比較する実験により，
 誤りや確認回数の減少，理解の容易さ・自然さ・協調効率の向上を示した．
 ゆえに，ロボットに共同注意を実装すると，意図の可視化と曖昧性の低減を通じて，
 人とのより自然で効率的なコミュニケーションが可能になる．[6]
 さらに，移動ロボットが人間と共同注意を成立させるには，相手の相互作用空間を考慮し，適切な位置へ移動する必要がある．
\subsection{F-formationに基づく共同注意の確立}
 近接学の第一人者 Kendon は，2 人以上が対面して相互行為を行う際に観察される空間的配置を F 陣形と定義し
 その枠組みは社会工学の基盤理論として広く用いられている．[7] 本セクションでは人物の位置と向きから会話グループとF陣形を生成する方法について、
 調査を行った.
 
 Setti らは，単眼カメラで得られる人物の位置と向きだけを使い，
 GCFF と呼ばれる手法で会話グループの特定と F 陣形の検出を実現した．[8]
 考え方はシンプルで，まず人の立ち位置と向きから，参加者全員が等しくアクセスできる共通の中心領域（o-space）の候補を作り，その割り当てを最も整合的にするよう最適化する．
 さらに，可視性の制約を入れることで，誰かが遮られて中心を見通せない配置はグループとして認めない．
 実験では，雑音や人数の違いがあっても安定してグループと o-space を見つけられることが示され，位置と向きという基本情報だけで社会的なやり取りの場を十分に捉えられることが確認された．
 これにより，ロボットは人の集まりの構造を理解し，どこに立てば共同注意を保てるかを判断するための土台を得られる．
 
 Barua らは，ロボット側のカメラ映像から人の骨格キーポイントを抽出し，
 それにもとづいて F 陣形の種類とロボットが近づくべき角度を同時に推定する手法を提案した.[9]
 処理はリアルタイムで動き，その場の人が会話の輪に入っているかどうかを見分ける．
 そして，F陣形を検出しながら，ロボットがどの方向から近づき，どこで止まればよいかを決める．
 具体的には，カメラで得た人の位置と向きを学習モデルに入力し，どちら側から入ると自然か，
 どの向きで止まれば相手の視線をさえぎらず，声も届きやすいかを判断する．
 実験では，グループ検出や接近角度の推定で高い精度を示し，従来法よりも安定して社会的に受け入れられる立ち位置を選べた．
 結果として，ロボットは人のやり取りに無理なく参加し，共同注意を保ちながら，より自然で予測しやすい対話へとつなげられることが示された．
 
\section{対話を使った自律移動}
 HRIの視点では，ロボットが人の意図を理解し，配慮しながら動くことが重要である．
 LLM を導入すると，あいまいな言い方や文脈も読み取り，必要に応じて質問して確認し，人に合わせて行動を調整できる．
 これにより，安全・わかりやすさといった人に中心の考え方に沿って，距離の取り方や見せ方，合図などを対話で決められる．
 さらに，目的地案内だけでなく，情報の提示や物の確認，複数手順の依頼など，より広く複雑なタスクも一貫して対応できる．
 本セクションは，対話を用いてロボットが自律的に目的地へ移動できるようにする方法を調査する．

 AhnらはLLM が理解した指示をそのまま実行させるのではなく，
 「いまのロボットに本当にできるか」を合わせて判断し，
 次にやる行動を一つずつ選んで実行する方法を提案した.
 まず LLM が指示に対して「この技能がふさわしいか」をスコア化し，
 次にロボット側の学習済み価値関数や成功確率から
 「この技能は今の状態で実行できそうか」をスコア化する．
 最後に，この二つを掛け合わせて一番よい技能を選び，
 実行してはロボットの状態を更新し，また次の行動を選ぶ――という流れを繰り返す．
 この枠組みにより，
 LLM の柔らかな理解力（あいまいな言い方の解釈，手順の分解）を，
 センサー・移動能力など現実の制約に接地できる．
 実験では，複数ステップの家事タスクなどで，
 間違った行動を減らし，長い指示を安全に分解して進められることが示された．
 一方で，あらかじめ用意した技能や価値関数の質に性能が依存し，
 未登録の行動や環境の大きな変化には弱いという課題もある．
 ロボットにとっては，「言語でわかる」だけでなく
 「今できることを選ぶ」仕組みが手に入り，
 対話に基づく自律的な計画と実行の土台になる．
 
 
 \section{おわりに}



% % 参考文献：bibtexを使う場合
% \bibliographystyle{junsrt} %参考文献出力スタイル .bst
% \bibliography{mabst} %bibtexファイル .bib
% \label{sannkoubunnkenn_chapter}


%% 参考文献：手書きの場合
\begin{thebibliography}{99}
\label{sannkoubunnkenn_chapter}

%% 程研究室参考文献形式
% 学術雑誌論文
\bibitem{ChengXX}
ほげほげ: あべば入門, 何らか出版社, 200?年?月.

\end{thebibliography}


\end{document}
